\documentclass[12pt]{article}
\usepackage{amsmath, amsthm, amssymb}
\usepackage{geometry}
\geometry{margin=1in}
\usepackage{hyperref}
\usepackage{booktabs}

\newtheorem{theorem}{Theorem}
\newtheorem{proposition}[theorem]{Proposition}
\newtheorem{corollary}[theorem]{Corollary}
\theoremstyle{definition}
\newtheorem{definition}[theorem]{Definition}
\theoremstyle{remark}
\newtheorem{remark}[theorem]{Remark}

\title{Mathematics as Cohomological Grammar}
\author{Diego Rinc\'on\\
\small Independent}
\date{}

\begin{document}

\maketitle

\begin{abstract}
A grammar specifies what combinations are valid.
Cohomology specifies what combinations persist ---
what survives every valid transformation.
We propose that mathematics is the study of cohomological grammars:
not the rules of combination, but their invariants.
A theorem is a cohomology class.
A proof is a homotopy.
A conjecture is a closed form whose exactness is unknown.
The plurality of mathematical theories corresponds to the plurality
of cohomology theories, each reading a different kind of persistence
from the same underlying structure.
Motivic cohomology is the universal grammar from which all others derive.
\end{abstract}

\section{Grammar and Persistence}

A grammar has two parts: what can be combined (syntax) and what the
combination means (semantics). Classical grammar asks whether a sequence
is well-formed. It does not ask whether the sequence survives.

Cohomology asks only the second question.

\begin{definition}[Cohomological grammar]
A \emph{cohomological grammar} on a space $X$ is a cohomology theory
$H^*$ together with:
\begin{enumerate}
\item a \emph{combination rule}: the cup product
  $\smile: H^p \otimes H^q \to H^{p+q}$,
\item a \emph{validity condition}: the cocycle condition $d\omega = 0$
  (the combination is internally consistent),
\item a \emph{triviality condition}: the coboundary condition
  $\omega = d\eta$ (the combination is explained by something simpler).
\end{enumerate}
The grammar is the data $(H^*, \smile, d)$.
A class $[\omega] \in H^p(X)$ is a \emph{valid combination that is not trivial}.
\end{definition}

\begin{remark}
The syntax/semantics distinction collapses here.
A cocycle is both a valid expression (satisfies the rule $d\omega = 0$)
and a meaning (the class $[\omega]$ it represents in $H^p$).
Validity and content are the same thing.
\end{remark}

\section{Mathematical Objects as Classes}

\begin{definition}[Theorem, proof, conjecture]
In a cohomological grammar on a mathematical space $\mathcal{M}$:
\begin{itemize}
\item A \emph{theorem} is a cohomology class $[\omega] \in H^p(\mathcal{M})$:
  a combination that is valid ($d\omega = 0$) and non-trivial
  ($\omega \neq d\eta$ for any $\eta$).
\item A \emph{proof} is a homotopy: a continuous path in the space
  of valid combinations from one representation of $[\omega]$ to another.
  Two proofs are equivalent if they represent the same class.
\item A \emph{conjecture} is a closed form $\omega$ with $d\omega = 0$
  whose class $[\omega]$ in $H^*$ is unknown: it might be zero
  (provably follows from known structure) or non-trivial (genuinely new).
\item An \emph{open problem} is the question: is $[\omega] = 0$?
\end{itemize}
\end{definition}

\begin{proposition}[The hard problems are non-trivial classes]
The Clay Millennium Problems \cite{hard} are closed forms in the
motivic cohomology of $\mathrm{Spec}\,\mathbb{Z}$ whose triviality
is unknown.
Each has been verified locally at every prime
(every $\ell$-adic and $p$-adic incarnation is trivial),
but the global class has not been shown to vanish.
This is the arithmetic statement of the Arithmetic Wall \cite{synthesis}:
locally trivial, globally unresolved.
\end{proposition}

\section{The Plurality of Grammars}

Different cohomology theories are different grammars applied to
the same space. They read different kinds of persistence.

\begin{definition}[Grammar dictionary]
\hfill
\begin{center}
\begin{tabular}{@{}lll@{}}
\toprule
Grammar & Persistence type & What it reads \\
\midrule
de Rham $H^*_{\rm dR}$ & Smooth deformation & Differential forms \\
Singular $H^*_{\rm sing}$ & Continuous deformation & Cycles and boundaries \\
\'Etale $H^*_{\rm\acute{e}t}$ & Algebraic deformation & Galois action \\
Crystalline $H^*_{\rm crys}$ & $p$-adic deformation & Frobenius structure \\
Motivic $H^*_{\rm mot}$ & All deformations & Universal \\
\bottomrule
\end{tabular}
\end{center}
Each grammar is a different answer to the question:
\emph{what does it mean for a combination to persist?}
\end{definition}

\begin{theorem}[Motivic grammar is universal \cite{motives}]
For a smooth variety $X$ over a field $k$,
motivic cohomology $H^*_{\rm mot}(X, \mathbb{Z}(*))$ maps to every
other cohomology theory:
\[
H^*_{\rm mot}(X) \longrightarrow H^*_{\rm dR}(X),\quad
H^*_{\rm mot}(X) \longrightarrow H^*_{\rm\acute{e}t}(X),\quad
H^*_{\rm mot}(X) \longrightarrow H^*_{\rm crys}(X),
\]
via comparison isomorphisms (over $\mathbb{C}$) or regulator maps
(over number fields). Every grammar is a quotient of the motivic grammar.
A statement true in the motivic grammar is true in every grammar.
\end{theorem}

\begin{remark}[Mathematics as the field of grammars]
Mathematics is not one grammar but the field in which grammars
discover each other.
The comparison isomorphisms --- de Rham equals singular over $\mathbb{C}$,
étale equals crystalline via $B_{\rm crys}$ \cite{perfectoid} ---
are the moments when two grammars, developed independently,
turn out to be reading the same persistence from the same space.
These moments are not coincidences. They are the content of mathematics.
\end{remark}

\section{Proof as Homotopy}

The identification of proofs with homotopies is not a metaphor.
In homotopy type theory (Voevodsky \cite{voevodsky2013}),
a proof of a proposition $P$ is a term of type $P$.
Two proofs are equal if there is a higher proof (a homotopy) between them.
A theorem is a type that is inhabited; its cohomology class
is the connected component of the space of proofs.

\begin{proposition}[Proof irrelevance and coboundaries]
In a cohomological grammar, two proofs $\eta_1, \eta_2$ of the same
theorem $[\omega]$ satisfy $\omega = d\eta_1 = d\eta_2$.
They differ by $d(\eta_1 - \eta_2) = 0$, i.e., $\eta_1 - \eta_2$ is a
closed form of one degree lower.
\emph{Proof irrelevance} is the statement that the choice of $\eta$
does not affect $[\omega]$: the theorem is the class, not the cocycle.
\end{proposition}

\begin{remark}[Long proofs and higher cohomology]
A very long proof is a high-dimensional coboundary: many intermediate
steps, each a low-dimensional form, assembled into a single boundary.
A short proof is a low-dimensional coboundary of the theorem.
The length of a proof is related to the dimension of the ``filling''
required to show the theorem is trivial in a higher grammar.

This is why the Millennium Problems have long proofs in low-dimensional
grammars but no proof at all in the motivic grammar: the filling
requires a higher-dimensional object (a Frobenius, an Euler system,
a global section) that the lower grammars cannot construct.
\end{remark}

\section{The Grammar of the Arithmetic Wall}

\begin{proposition}[The Wall as a grammatical boundary]
The Arithmetic Wall \cite{synthesis} is the boundary between two grammars:
\begin{enumerate}
\item The \emph{local grammar}: étale cohomology at each prime $p$,
  equipped with the Frobenius $\varphi_p$. This grammar is complete:
  every class is either trivial or has a known non-trivial certificate
  (an eigenvalue of Frobenius).
\item The \emph{global grammar}: motivic cohomology over $\mathbb{Q}$,
  with no Frobenius. This grammar is incomplete:
  closed forms exist whose triviality is unknown.
\end{enumerate}
The Wall is not a theorem and not a conjecture.
It is the statement that the local grammar does not generate the global one ---
that local persistence does not imply global persistence.
This is a grammatical fact, not an arithmetical one.
\end{proposition}

\begin{remark}[What a proof of BSD would be]
A proof of BSD would be a coboundary in the motivic grammar:
a form $\eta$ such that
$d\eta = \omega_{\rm BSD}$,
where $\omega_{\rm BSD}$ is the closed form encoding the rank-$L$-value
comparison.
The Iwasawa main conjecture \cite{iwasawa} constructs $\eta$ locally at $p$.
The Bloch--Kato conjecture \cite{blochkato} asserts $\eta$ exists globally.
Neither proves it.
The proof would be: exhibit the global filling.
\end{remark}

\section{Grammar and Incompleteness}

G\"odel's incompleteness theorem is a grammatical theorem.

A formal system is a grammar: axioms (valid base combinations)
and inference rules (valid extension operations).
A provable sentence is a coboundary: reachable from the axioms
by the inference rules, i.e., $\omega = d\eta$ for some derivation $\eta$.
An unprovable sentence is a closed form that is not a coboundary:
$d\omega = 0$ (it is internally consistent) but $\omega \neq d\eta$
for any derivation in the system.

G\"odel's sentence is a closed form that is not a coboundary
in the system's own grammar.

\begin{proposition}[The Millennium Problems as G\"odel sentences]
It is possible that the Millennium Problems are closed forms
that are not coboundaries in the current motivic grammar ---
not because the grammar is too weak to construct the filling,
but because the filling requires axioms not yet adopted.
A proof of RH may require extending the grammar of arithmetic
by a new axiom (the existence of a global Frobenius over $\mathbb{Q}$,
or a new cohomological structure) from which the proof proceeds.
The wall would then be not a missing argument
but a missing grammar rule.
\end{proposition}

\begin{remark}
This is consistent with the history.
The Weil conjectures were closed forms unprovable
in the grammar of classical algebraic geometry.
Grothendieck extended the grammar (\'etale cohomology, topoi, motives)
and Deligne proved them.
The extension was not a trick. It was a new grammar.
The proof was a coboundary in the new grammar
that was not a coboundary in the old one.
\end{remark}

\section{What a New Grammar Looks Like}

A new grammar does not arrive as a proof.
It arrives as a definition.

Grothendieck did not prove the Weil conjectures.
He defined a new kind of object (a scheme)
and a new kind of covering (an \'etale site)
and a new kind of cohomology (\'etale cohomology)
from which Deligne could build the coboundary.
The new grammar had to be built first.
The proof was the last step, not the hard one.
The hard step was the grammar.

\begin{remark}[What the proof of RH will be]
The proof of RH will not be a clever manipulation
of the existing grammar.
It will be a new grammar
in which the closed form $\omega_{\rm RH}$ becomes a coboundary:
in which the Hadamard product control becomes a
computable filling of dimension one higher.
What candidates exist for the new grammar?
The Fargues--Fontaine curve \cite{perfectoid}:
the geometric bridge between characteristic $0$ and $p$.
Connes' spectral triple on $\mathbf{A}_\mathbf{Q}/\mathbf{Q}^*$ \cite{connes}:
where the absorption frequencies are the zeros.
A global Frobenius over $\mathrm{Spec}\,\mathbb{Z}$:
the missing symmetry that the local grammars have and the global one does not.
None of these is the new grammar yet.
Each is a face of it.
\end{remark}

\section{One Grammar or Many}

The title carries a parenthetical: grammar(s).
The question is whether mathematics has one grammar or many.

\begin{theorem}[One and many]
Mathematics has one grammar (motivic cohomology) and many grammars
(all the rest). The many are quotients of the one.
The one is not accessible directly: it must be approached through the many.
This is the mathematical analogue of a language that can only be spoken
in translation.
\end{theorem}

\begin{remark}
The Langlands program \cite{langlands} is the grammar of translations:
it specifies which automorphic grammar (the many, over the adeles)
corresponds to which Galois grammar (the many, over $\bar{\mathbb{Q}}$).
The motivic grammar (the one) is what they are both translating.
The Langlands correspondence is a dictionary.
The motive is the word.
\end{remark}

\begin{remark}[Grammar and thought]
A mathematical grammar is a grammar of thought \cite{grammars}.
It specifies which thoughts are thinkable ---
which combinations are valid, which persist under transformation.
The motivic grammar is the deepest known grammar of what persists in mathematics.
It is not the deepest grammar of thought.
Thought occurs in timespace \cite{timespace}:
the grammar runs, the running takes time,
the thought is the running.
The motivic grammar is what survives when the running stops.
It is the cohomology of mathematical thought:
the classes that persist after all the proofs have been forgotten.
\end{remark}

\begin{thebibliography}{9}

\bibitem{hard}
Rinc\'on, D. (2026).
The Hard Step: Why the Clay Millennium Problems Resisted Solution. Preprint.

\bibitem{synthesis}
Rinc\'on, D. (2026).
The Arithmetic Wall. Preprint.

\bibitem{motives}
Rinc\'on, D. (2026).
Motives, Standard Conjectures, and the Unification of $L$-functions. Preprint.

\bibitem{perfectoid}
Rinc\'on, D. (2026).
Perfectoid Spaces and the Local Frobenius. Preprint.

\bibitem{langlands}
Rinc\'on, D. (2026).
The Langlands Thread. Preprint.

\bibitem{iwasawa}
Rinc\'on, D. (2026).
Iwasawa Theory: The $p$-adic Machine for BSD and Beyond. Preprint.

\bibitem{blochkato}
Rinc\'on, D. (2026).
The Bloch--Kato Conjecture. Preprint.

\bibitem{voevodsky2013}
Voevodsky, V.\ et al.\ (2013).
\textit{Homotopy Type Theory: Univalent Foundations of Mathematics}.
Institute for Advanced Study.

\bibitem{one_wall}
Rinc\'on, D. (2026).
One Wall: The Local-to-Global Conjecture. Preprint.

\bibitem{grammars}
Rinc\'on, D. (2026).
Grammar(s): Mind(s) versus Mentality: Language(s) of Being(s). Preprint.

\bibitem{timespace}
Rinc\'on, D. (2026).
Thought Is Language in Timespace. Preprint.

\end{thebibliography}

\end{document}
